% Standalone preview of the illustrated FC graph-construction pipeline figure.
% Mirrors THESIS/settings.tex's package load order, TUM colour definitions, and
% light/dark theme switch so the figure previews exactly as it would in the
% thesis, without touching THESIS/ itself. Compile:
%   tectonic -o build preview.tex        (light theme)
%   tectonic -o build preview_dark.tex   (dark theme, via \jobname switch)

\documentclass[fontsize=11pt,paper=a4]{scrbook}

\usepackage[utf8]{inputenc}
\usepackage[T1]{fontenc}
\usepackage{amsmath}
\usepackage{amssymb}
\usepackage[sc]{mathpazo}
\usepackage{graphicx}
\graphicspath{{./}}
\usepackage{array}
\usepackage{tikz}
\usetikzlibrary{arrows.meta,positioning,calc}
\usepackage{caption}
\usepackage{xcolor}
\usepackage[printonlyused]{acronym}
\usepackage{ifthen}

% Light/dark theme switch, identical mechanism to THESIS/settings.tex.
\ifthenelse{\equal{\detokenize{dark}}{\jobname}}{%
  \newcommand{\bg}{black}
  \newcommand{\fg}{white}
  \usepackage[pagecolor=\bg]{pagecolor}
  \color{\fg}
}{%
  \newcommand{\bg}{white}
  \newcommand{\fg}{black}
}

% TUM corporate design colours, copied from THESIS/settings.tex.
\definecolor{TUMBlue}{HTML}{0065BD}
\definecolor{TUMSecondaryBlue}{HTML}{005293}
\definecolor{TUMSecondaryBlue2}{HTML}{003359}
\definecolor{TUMBlack}{HTML}{000000}
\definecolor{TUMWhite}{HTML}{FFFFFF}
\definecolor{TUMDarkGray}{HTML}{333333}
\definecolor{TUMGray}{HTML}{808080}
\definecolor{TUMLightGray}{HTML}{CCCCC6}
\definecolor{TUMAccentGray}{HTML}{DAD7CB}
\definecolor{TUMAccentOrange}{HTML}{E37222}
\definecolor{TUMAccentGreen}{HTML}{A2AD00}
\definecolor{TUMAccentLightBlue}{HTML}{98C6EA}
\definecolor{TUMAccentBlue}{HTML}{64A0C8}

\begin{document}

\pagecolor{\bg}
\color{\fg}

\begin{acronym}
  \acro{ROI}[ROI]{region of interest}
  \acro{MCI}[MCI]{mild cognitive impairment}
  \acro{BOLD}[BOLD]{blood-oxygen-level-dependent}
  \acro{fMRI}[fMRI]{functional magnetic resonance imaging}
\end{acronym}

\noindent\textbf{Preview page geometry:} \the\textwidth\ (compare against the thesis
\texttt{scrbook} textwidth before sizing \texttt{\textbackslash resizebox} in
\texttt{THESIS/chapters/04\_methodology.tex}).

\begin{figure}[p]
  \centering
  \resizebox{\textwidth}{!}{% TikZ diagram: Step-by-Step Functional Connectivity Graph Construction Pipeline
% Illustrated version -- each stage card pairs a small rendered inset panel with
% shortened text. Lengthwise (Vertical) Flowchart for Chapter 4 (Section 4.2).
% Compatible with both Light and Dark thesis themes.
%
% Assumes graphicx, tikz (with libraries arrows.meta, positioning, calc) and the
% array package's m{} column type are already loaded by the enclosing document
% (THESIS/settings.tex loads all of these; see preview.tex for a standalone load).
% \includegraphics uses bare basenames so \graphicspath resolves them both here
% (graphicspath={{./}}) and inside THESIS (graphicspath={{figures/}}).

\begin{tikzpicture}[
    font=\small\sffamily,
    >=Latex,
    every node/.style={align=center},
    card1/.style={
        draw=TUMSecondaryBlue2!90,
        fill=white,
        rounded corners=5pt,
        line width=0.85pt,
        inner sep=5pt,
        font=\footnotesize,
        text=black!90
    },
    card2/.style={
        draw=TUMBlue!90,
        fill=TUMBlue!10,
        rounded corners=5pt,
        line width=0.85pt,
        inner sep=5pt,
        font=\footnotesize,
        text=black!90
    },
    card3/.style={
        draw=TUMAccentOrange!90,
        fill=TUMAccentOrange!12,
        rounded corners=5pt,
        line width=0.85pt,
        inner sep=5pt,
        font=\footnotesize,
        text=black!90
    },
    card4a/.style={
        draw=TUMBlue!90,
        fill=TUMBlue!10,
        rounded corners=5pt,
        line width=0.85pt,
        inner sep=5pt,
        font=\footnotesize,
        text=black!90
    },
    card4b/.style={
        draw=TUMAccentGreen!90,
        fill=TUMAccentGreen!15,
        rounded corners=5pt,
        line width=0.85pt,
        inner sep=5pt,
        font=\footnotesize,
        text=black!90
    },
    card5/.style={
        draw=TUMSecondaryBlue2!90,
        fill=TUMSecondaryBlue2!12,
        rounded corners=5pt,
        line width=0.95pt,
        inner sep=5pt,
        font=\footnotesize,
        text=black!90
    },
    card6/.style={
        draw=TUMBlue!90,
        fill=TUMBlue!12,
        rounded corners=5pt,
        line width=0.85pt,
        inner sep=5pt,
        font=\footnotesize,
        text=black!90
    },
    flowline/.style={
        ->,
        line width=1.0pt,
        draw=\fg!85,
        shorten >=1.5pt,
        shorten <=1.5pt
    },
    flowlabel/.style={
        font=\scriptsize\bfseries,
        text=\fg!95
    }
]

% Thin neutral-grey border around each embedded inset panel -- no fill, so
% the panel's transparent background lets the card tint show through.
\newcommand{\insetimg}[2]{\tikz[baseline]{\node[draw=TUMGray!60, line width=0.4pt, inner sep=0pt] {\includegraphics[width=#1]{#2}};}}

% -------------------------------------------------------------
% STAGE 1: 4D fMRI ACQUISITION & SCHAEFER-200 PARCELLATION
% -------------------------------------------------------------
\node[card1, text width=13.5cm] (card1) at (0, 0) {
    \begin{tabular}{@{}m{2.6cm}@{\hspace{7pt}}m{9.9cm}@{}}
        \insetimg{2.5cm}{fmri_input} &
        \textbf{1. 4D Resting-State BOLD fMRI Acquisition \& Atlas Parcellation}\newline
        \textnormal{\footnotesize Preprocessed functional volumes aligned to standard MNI152 space, parcellated into $|V| = 200$ cortical \acp{ROI} using the Schaefer-200 atlas (7 canonical functional networks).}
    \end{tabular}
};

% -------------------------------------------------------------
% STAGE 2: BOLD TIME SERIES EXTRACTION
% -------------------------------------------------------------
\node[card2, text width=13.5cm, below=0.7cm of card1] (card2) {
    \begin{tabular}{@{}m{2.6cm}@{\hspace{7pt}}m{9.9cm}@{}}
        \insetimg{2.5cm}{regional_timeseries} &
        \textbf{2. Regional BOLD Signal Averaging}\newline
        \textnormal{\footnotesize Mean timecourse per parcel $j \in \{1, \dots, 200\}$: $\displaystyle \mathbf{T} = \begin{bmatrix} \mathbf{t}_1^\top \\ \vdots \\ \mathbf{t}_{200}^\top \end{bmatrix} \in \mathbb{R}^{200 \times L}$, where $L$ is the number of temporal frames per scan.}
    \end{tabular}
};
\draw[flowline] (card1.south) -- (card2.north) node[midway, right=4pt, flowlabel] {Voxel-wise spatial averaging};

% -------------------------------------------------------------
% STAGE 3: FUNCTIONAL CONNECTIVITY & FISHER-Z
% -------------------------------------------------------------
\node[card3, text width=13.5cm, below=0.7cm of card2] (card3) {
    \begin{tabular}{@{}m{2.6cm}@{\hspace{7pt}}m{9.9cm}@{}}
        \insetimg{2.5cm}{fc_matrix} &
        \textbf{3. Pairwise Functional Connectivity Matrix \& Fisher $z$-Transformation}\newline
        \textnormal{\footnotesize Pearson correlation $r_{ij} = \operatorname{corr}(\mathbf{t}_i, \mathbf{t}_j)$, Fisher-transformed: $\displaystyle z_{ij} = \operatorname{arctanh}(r_{ij}) = \frac{1}{2}\ln\left(\frac{1+r_{ij}}{1-r_{ij}}\right)$, yielding symmetric $C \in \mathbb{R}^{200 \times 200}$ (zero diagonal).}
    \end{tabular}
};
\draw[flowline] (card2.south) -- (card3.north) node[midway, right=4pt, flowlabel] {Pairwise temporal coupling};

% -------------------------------------------------------------
% STAGE 4: DUAL GRAPH REPRESENTATION (SPLIT BRANCHES)
% -------------------------------------------------------------
% Branch 4a (Left): Node Features
\node[card4a, text width=6.3cm, below left=1.0cm and 0.3cm of card3.south] (card4a) {
    \begin{tabular}{@{}m{2.2cm}@{\hspace{5pt}}m{3.6cm}@{}}
        \insetimg{2.1cm}{node_features} &
        \textbf{4a. Regional Node Features}\newline
        \textnormal{\footnotesize Each region carries its full connectivity row: $\mathbf{x}_j = C_{j,:} \in \mathbb{R}^{200}$, so $X = C \in \mathbb{R}^{200 \times 200}$.}
    \end{tabular}
};

% Branch 4b (Right): Topological Adjacency
\node[card4b, text width=6.3cm, below right=1.0cm and 0.3cm of card3.south] (card4b) {
    \begin{tabular}{@{}m{2.2cm}@{\hspace{5pt}}m{3.6cm}@{}}
        \insetimg{2.1cm}{sparse_graph} &
        \textbf{4b. Symmetrised $k$-NN Adjacency}\newline
        \textnormal{\footnotesize Symmetrised top-$k$ over $|C_{ij}|$: $\displaystyle A = \max(A_{\text{top-}k}, A_{\text{top-}k}^\top) \in \{0,1\}^{200 \times 200}$.\newline \scriptsize $k=16$ (pretraining) $\,\big|\,$ $k=8$ (supervised)}
    \end{tabular}
};

\draw[flowline] (card4a.north |- card3.south) -- (card4a.north) node[midway, left=3pt, flowlabel] {Raw FC rows};
\draw[flowline] (card4b.north |- card3.south) -- (card4b.north) node[midway, right=3pt, flowlabel] {Topological sparsification};

% -------------------------------------------------------------
% STAGE 5: TIME-AUGMENTED LONGITUDINAL GRAPH SERIES
% -------------------------------------------------------------
\node[card5, text width=13.5cm, anchor=north] (card5) at ($(card3.south |- card4a.south)+(0, -1.0)$) {
    \begin{tabular}{@{}m{3.3cm}@{\hspace{7pt}}m{9.2cm}@{}}
        \begin{tabular}{@{}c@{\hspace{2pt}}c@{\hspace{2pt}}c@{}}
            \includegraphics[width=1.05cm]{graph_snapshot1} &
            \raisebox{0.35cm}{\tikz[baseline]{\draw[->, line width=0.8pt, draw=\fg!85] (0,0) -- (0.45,0) node[midway, above, font=\scriptsize] {$\Delta t_1$};}} &
            \includegraphics[width=1.05cm]{graph_snapshot2} \\[-1pt]
            \multicolumn{3}{c}{} \\[-9pt]
            & \raisebox{0.35cm}{\tikz[baseline]{\draw[->, line width=0.8pt, draw=\fg!85] (0,0) -- (0.45,0) node[midway, above, font=\scriptsize] {$\Delta t_2$};}} & \includegraphics[width=1.05cm]{graph_snapshot3}
        \end{tabular}
        &
        \textbf{5. Time-Augmented Longitudinal Graph Sequence}\newline
        \textnormal{\footnotesize For subject $i$ with $T_i$ visits, visit $t$ yields $\mathcal{G}_i^{(t)} = (V, A_i^{(t)}, X_i^{(t)})$ with $|V|=200$. Time cadence: normalized interval $\Delta t_t = \frac{\text{cum}_t - \text{cum}_{t-1}}{108}$ or cumulative $\log(1+\text{months})$.\newline $\mathcal{S}_i = \bigl( (\mathcal{G}_i^{(1)}, \Delta t_1), \dots, (\mathcal{G}_i^{(T_i)}, \Delta t_{T_i}) \bigr)$}
    \end{tabular}
};

\draw[flowline] (card4a.south) -- (card4a.south |- card5.north);
\draw[flowline] (card4b.south) -- (card4b.south |- card5.north);

% -------------------------------------------------------------
% STAGE 6: DOWNSTREAM SUBJECT-LEVEL TARGET
% -------------------------------------------------------------
\node[card6, text width=13.5cm, below=0.7cm of card5] (card6) {
    \begin{tabular}{@{}m{2.6cm}@{\hspace{7pt}}m{9.9cm}@{}}
        \insetimg{2.5cm}{subject_outcome} &
        \textbf{6. Downstream Subject-Level Conversion Target}\newline
        \textnormal{\footnotesize Single binary label per patient trajectory: $y_i \in \{0,1\}$ ($y_i=0$ for Stable \ac{MCI}, $y_i=1$ for Converter \ac{MCI}).\newline \scriptsize Effective sample size is bounded by subject count $N$ (DELCODE $n=167$, Pooled $n=248$), not scan count.}
    \end{tabular}
};
\draw[flowline] (card5.south) -- (card6.north) node[midway, right=4pt, flowlabel] {Sequence-to-label mapping};

\end{tikzpicture}
}
  \caption{Subject-level functional connectivity graph construction pipeline
    (illustrated, preview). (1) Resting-state four-dimensional \ac{BOLD} \ac{fMRI}
    scans are parcellated into $|V| = 200$ cortical \acp{ROI} using the Schaefer-200
    atlas. (2) Mean BOLD time series ($\mathbf{T} \in \mathbb{R}^{200 \times L}$) are
    extracted per region. (3) Pairwise Pearson correlation followed by Fisher
    $z$-transformation yields the continuous functional connectivity matrix
    $C \in \mathbb{R}^{200 \times 200}$. (4a) Raw connectivity rows form 200-dimensional
    node feature vectors ($X = C$). (4b) Symmetrised $k$-nearest-neighbour
    sparsification over $|C|$ constructs the binary adjacency matrix $A$. (5) The
    resulting visit graph $\mathcal{G}_i^{(t)}$ is paired with normalized interval
    timing $[\Delta t_t]$ across the longitudinal sequence. Inset panels (brain
    slices, connectivity heatmaps, node/edge diagrams) are synthetic schematic
    illustrations of the construction steps, not measured subject data.}
  \label{fig:graph-construction-pipeline-preview}
\end{figure}

\end{document}
